%--------------------------------------
%  Modelo com instruções para preparação do resumo para submissão de trabalho no VIII Simpósio de Matemática da Região Sul Fluminense 2022
% Autor: Ivan Aguilar - Departamento de Matemática - ICEx-UFF
%--------------------------------------

\documentclass[a4paper,10pt]{article}

\usepackage[brazil]{babel}
\usepackage[T1]{fontenc}
\usepackage[utf8]{inputenc}

\usepackage{amsfonts, amsmath, amssymb, amsthm}
\usepackage[pdftex]{graphicx}
\usepackage{epstopdf}
\usepackage{fancyhdr}
\usepackage{epsfig}
\usepackage{xcolor}
\usepackage{hyperref}
\hypersetup{
    colorlinks,
    linkcolor={red!50!black},
    citecolor={blue!50!black},
    urlcolor={blue!80!black}
}
\urlstyle{same}

\newcommand{\headrulecolor}[1]{\patchcmd{\headrule}{\hrule}{\color{#1}\hrule}{}{}}
\renewcommand{\headrulewidth}{0.5pt}
\fancyhead[R]{}
\fancyhead[C]{\begin{minipage}{4cm}
\includegraphics[width=4cm]{logosimma.pdf}
\end{minipage}
\begin{minipage}{11.5cm}
    \begin{flushright}
       \footnotesize{VIII Simpósio de Matemática da Região Sul Fluminense -- SIMMA 2022\\ Universidade Federal Fluminense -- UFF, Volta  Redonda, RJ}
    \end{flushright}
\end{minipage}
}

\pagenumbering{gobble}
\headheight 20mm      %
\oddsidemargin 2.0mm  %
\evensidemargin 2.0mm %
\topmargin -25mm      %
\textheight 250mm     %
\textwidth 160mm      %
%---- FIM DO PREÂMBULO --------%

%---- INÍCIO DO DOCUMENTO --------%

\begin{document}
\title{\vspace{-1cm} 
\Large{\bf Instruções para Preparação do Resumo a ser Submetido\\ 
ao VIII Simpósio de Matemática da Região Sul Fluminense}
}

\author{ 
{\bf {\large Pitágoras de Samos}} \hspace*{1cm} 
{\bf {\large \underline{Hypatia de Alexandria}\footnote{Bolsista de Iniciação Científica PIBIC/CNPq}}}  \\[4pt]
{\small Instituto Pitagórico, UFF} \\
{\small Av. Platão da Grecia, 314} \\
{\small 12345-000, Volta Redonda, RJ} \\
{\small E-mail: pitagoras@id.uff.br \hspace*{1cm}  hypatia@id.uff.br} 
\\
\\
{\bf {\large Henry Poincaré}}  \\[4pt]
{\small Instituto de Matemática, UFF} \\
{\small Rua Fermat, 35} \\
{\small 1000-100, Planolândia, PI} \\
{\small E-mail: plano@id.uff.br}
 \\
 \\
 \\
\textbf{RESUMO}
}

\date{}
\maketitle
\thispagestyle{fancy}

%------ INÍCIO DO RESUMO ---------%

Os trabalhos a serem submetidos para o VIII Simpósio de Matemática da Região Sul Fluminense, devem ter uma ou duas páginas, sem numeração. O nome do apresentador do trabalho deve ser sublinhado. Estas instruções foram preparadas no formato em que o resumo deve ser apresentado. 


As principais referências bibliográficas devem estar em ordem alfabética pelo sobrenome do primeiro autor. Como exemplo, a referência bibliográfica \cite{hig} refere-se a livro e a referência
\cite{key1}, a artigo em revista.

O arquivo com o resumo deve ser enviado em formato \textbf{PDF} a \href{mailto:simma.vma.vcx@id.uff.br}{\tt\color{blue}simma.vma.vcx@id.uff.br}.  O prazo final para submissão dos resumos é \textbf{\color{red} 15 de setembro de 2022}. O resultado da seleção de trabalhos aceitos será divulgado até o dia 16 de setembro de 2022 na página do SIMMA.

Mais informações em: {\color{blue}\url{http://simma.uff.br}}


%----- REFERÊNCIAS --------%
\begin{thebibliography}{99}

\bibitem{key1} BURGER, M.; HACKL, B.; RING, W. Incorporating topological derivatives
into level set methods. \emph{Journal of Computational Physics}, v. 194, n. 1, p. 344-362,
2004.

\bibitem{hig}  LITTLE, R. W. \emph{Elasticity}. New Jersey: Prentice-Hall, 1973.

\end{thebibliography}

\end{document}
